\documentclass[11pt,a4paper]{article}
\usepackage{get-uebung}

\blatt{9}{Überlagerungssatz}{Kap.~2.10}

\begin{document}

\blattkopf{
  \item Quellen korrekt deaktivieren (Spannungsquelle $\to$ Kurzschluss, Stromquelle $\to$ Leerlauf),
  \item Ströme und Spannungen in Netzwerken mit mehreren Quellen durch Überlagerung von Teilergebnissen berechnen,
  \item die Vorzeichen der Teilergebnisse mit Hilfe der Zählpfeile sicher bestimmen,
  \item begründen, warum Leistungen \emph{nicht} überlagert werden dürfen.
}

% -------------------------------------------------------------------
\begin{aufgabe}{Zwei Spannungsquellen}{15}
\begin{center}
\begin{circuitikz}
  \draw (0,0) to[V, v<=$U_1$] (0,3) to[R, l=$R_1$, i=$I_1$] (3,3)
        to[R, l=$R_2$, i=$I_2$] (6,3) to[V, v=$U_2$] (6,0) -- (0,0);
  \draw (3,3) to[R, l_=$R_3$, v^=$U_3$, i>_=$I_3$, *-*] (3,0);
\end{circuitikz}
\end{center}
\gegeben{$U_1 = \SI{12}{V}$, $U_2 = \SI{6}{V}$, $R_1 = \SI{4}{\ohm}$, $R_2 = \SI{2}{\ohm}$, $R_3 = \SI{4}{\ohm}$}
\begin{teilaufgaben}
  \item Zeichnen Sie die beiden Teilschaltungen, in denen jeweils nur eine Quelle aktiv ist.
  \item Berechnen Sie in beiden Teilschaltungen die Spannung $U_3$ (Teilergebnisse $U_3'$ und $U_3''$) und daraus $U_3$ und $I_3$.
  \item Bestimmen Sie auf dieselbe Weise die Ströme $I_1$ und $I_2$. Kontrollieren Sie Ihr Ergebnis mit dem Knotensatz.
  \item Berechnen Sie die Leistung $P_3$ in $R_3$. Vergleichen Sie mit der Summe $P_3' + P_3''$ der Leistungen aus den Teilschaltungen. Was folgt daraus?
\end{teilaufgaben}
\end{aufgabe}

\begin{loesung}
\schritt{a) Teilschaltungen}
\leitfrage{Wodurch wird eine deaktivierte \emph{Spannungs}quelle ersetzt? (Eselsbrücke: „Lass den Kreis der Quelle weg!“)}
Eine deaktivierte Spannungsquelle liefert $U = 0$, sie wird also zum Kurzschluss:
\begin{center}
\begin{circuitikz}[scale=0.75, transform shape]
  \draw (0,0) to[V, v<=$U_1$] (0,3) to[R, l=$R_1$, i=$I_1'$] (3,3)
        to[R, l=$R_2$, i=$I_2'$] (5,3) -- (5,0) -- (0,0);
  \draw (3,3) to[R, l_=$R_3$, v^=$U_3'$, *-*] (3,0);
  \node at (2.5,-0.7) {nur $U_1$ aktiv};
\end{circuitikz}
\hspace{1.2cm}
\begin{circuitikz}[scale=0.75, transform shape]
  \draw (0,0) -- (0,3) to[R, l=$R_1$, i=$I_1''$] (2,3)
        to[R, l=$R_2$, i=$I_2''$] (5,3) to[V, v=$U_2$] (5,0) -- (0,0);
  \draw (2,3) to[R, l_=$R_3$, v^=$U_3''$, *-*] (2,0);
  \node at (2.5,-0.7) {nur $U_2$ aktiv};
\end{circuitikz}
\end{center}
\schritt{b) Teilspannungen $U_3'$ und $U_3''$}
\leitfrage{Welche Widerstände liegen in der linken Teilschaltung parallel, welche in Reihe?}
Nur $U_1$: Der rechte Anschluss von $R_2$ ist über den Kurzschluss mit dem unteren Knoten verbunden, $R_2$ liegt also parallel zu $R_3$. Spannungsteiler aus $R_1$ und $R_2\parallel R_3 = \SI{4}{\ohm}\cdot\SI{2}{\ohm}/\SI{6}{\ohm} = \SI{1.33}{\ohm}$:
\[
  U_3' = U_1\cdot\frac{R_2\parallel R_3}{R_1 + R_2\parallel R_3} = \SI{12}{V}\cdot\frac{\SI{1.33}{\ohm}}{\SI{5.33}{\ohm}} = \SI{3}{V}
\]
Nur $U_2$: Jetzt liegt $R_1$ parallel zu $R_3$: $R_1\parallel R_3 = \SI{2}{\ohm}$.
\[
  U_3'' = U_2\cdot\frac{R_1\parallel R_3}{R_2 + R_1\parallel R_3} = \SI{6}{V}\cdot\frac{\SI{2}{\ohm}}{\SI{4}{\ohm}} = \SI{3}{V}
\]
Beide Quellen haben ihren Pluspol oben, beide Teilspannungen sind also positiv bezüglich des Pfeils $U_3$:
\[
  U_3 = U_3' + U_3'' = \ergebnis{\SI{6}{V}}, \qquad I_3 = \frac{U_3}{R_3} = \ergebnis{\SI{1.5}{A}}
\]
\schritt{c) Ströme $I_1$ und $I_2$}
Nur $U_1$: $I_1' = (U_1 - U_3')/R_1 = \SI{9}{V}/\SI{4}{\ohm} = \SI{2.25}{A}$; $I_2' = U_3'/R_2 = \SI{3}{V}/\SI{2}{\ohm} = \SI{1.5}{A}$.\\
Nur $U_2$: Der Strom fließt jetzt von rechts nach links, also \emph{gegen} die Pfeile $I_1$ und $I_2$:
\[
  I_1'' = -\frac{U_3''}{R_1} = -\frac{\SI{3}{V}}{\SI{4}{\ohm}} = -\SI{0.75}{A}, \qquad
  I_2'' = -\frac{U_2 - U_3''}{R_2} = -\frac{\SI{3}{V}}{\SI{2}{\ohm}} = -\SI{1.5}{A}
\]
Überlagerung:
\[
  I_1 = \SI{2.25}{A} - \SI{0.75}{A} = \ergebnis{\SI{1.5}{A}}, \qquad
  I_2 = \SI{1.5}{A} - \SI{1.5}{A} = \ergebnis{\SI{0}{A}}
\]
Knotensatz am oberen Knoten: $I_1 = I_2 + I_3 \Rightarrow \SI{1.5}{A} = \SI{0}{A} + \SI{1.5}{A}$ \checkmark\\
Anschaulich: $U_3 = \SI{6}{V} = U_2$ -- an $R_2$ liegt keine Spannung, also fließt kein Strom.
\schritt{d) Leistung -- die Falle}
\[
  P_3 = \frac{U_3^2}{R_3} = \frac{(\SI{6}{V})^2}{\SI{4}{\ohm}} = \ergebnis{\SI{9}{W}}, \qquad
  P_3' + P_3'' = \frac{(\SI{3}{V})^2}{\SI{4}{\ohm}} + \frac{(\SI{3}{V})^2}{\SI{4}{\ohm}} = \SI{4.5}{W} \neq P_3
\]
Der Überlagerungssatz gilt nur für Größen, die \emph{linear} von den Quellen abhängen (Ströme und Spannungen). Die Leistung hängt quadratisch ab: $(U_3' + U_3'')^2 \neq U_3'^2 + U_3''^2$. Leistungen werden deshalb immer erst \emph{nach} der Überlagerung aus dem Gesamtstrom bzw.\ der Gesamtspannung berechnet.
\typfehler{Die nicht aktive Spannungsquelle einfach mit ihren Leitungen entfernen (Leerlauf statt Kurzschluss). Und: Teilströme ohne Blick auf den Zählpfeil addieren -- prüfen Sie für jede Teilschaltung, ob der Strom in oder gegen die Pfeilrichtung fließt.}
\end{loesung}

% -------------------------------------------------------------------
\begin{aufgabe}{Spannungs- und Stromquelle gemischt}{18}
\begin{center}
\begin{circuitikz}
  \draw (0,0) to[V, v<=$U$] (0,3) to[R, l=$R_1$] (3,3)
        to[R, l=$R_3$, i=$I_3$] (6,3);
  \draw (3,3) to[R, l_=$R_2$, v^=$U_A$, *-*] (3,0);
  \draw (6,3) to[R, l_=$R_4$, v^=$U_B$, *-*] (6,0);
  \draw (8,0) to[I, i=$I_q$] (8,3) -- (6,3);
  \draw (0,0) -- (8,0);
  \node[above] at (3,3.35) {A};
  \node[above] at (6,3.35) {B};
\end{circuitikz}
\end{center}
\gegeben{$U = \SI{16}{V}$, $I_q = \SI{1}{A}$, $R_1 = R_2 = R_4 = \SI{20}{\ohm}$, $R_3 = \SI{10}{\ohm}$}
\begin{teilaufgaben}
  \item Wodurch ersetzen Sie die Stromquelle, wenn nur die Spannungsquelle wirkt? Zeichnen Sie beide Teilschaltungen.
  \item Berechnen Sie in der Teilschaltung „nur $U$“ die Spannungen $U_A'$, $U_B'$ und den Strom $I_3'$.
  \item Berechnen Sie in der Teilschaltung „nur $I_q$“ die Spannungen $U_A''$, $U_B''$ und den Strom $I_3''$.
  \item Geben Sie $U_A$, $U_B$ und $I_3$ an. In welche Richtung fließt der Strom durch $R_3$ tatsächlich?
  \item Welche Leistung gibt die Stromquelle ab?
\end{teilaufgaben}
\begin{tipp}Zeichnen Sie jede Teilschaltung sauber neu und suchen Sie nach Reihen- und Parallelschaltungen. Mit Spannungs- und Stromteiler kommen Sie dann ohne Gleichungssystem aus.\end{tipp}
\end{aufgabe}

\begin{loesung}
\schritt{a) Teilschaltungen}
Eine deaktivierte Stromquelle liefert $I = 0$ -- sie wird zum \emph{Leerlauf} (Zweig einfach weglassen). Die deaktivierte Spannungsquelle wird zum Kurzschluss:
\begin{center}
\begin{circuitikz}[scale=0.7, transform shape]
  \draw (0,0) to[V, v<=$U$] (0,3) to[R, l=$R_1$] (2.5,3) to[R, l=$R_3$, i=$I_3'$] (5,3);
  \draw (2.5,3) to[R, l_=$R_2$, v^=$U_A'$, *-*] (2.5,0);
  \draw (5,3) to[R, l_=$R_4$, v^=$U_B'$] (5,0) -- (0,0);
  \node at (2.5,-0.7) {nur $U$ aktiv ($I_q$ = Leerlauf)};
\end{circuitikz}
\hspace{0.8cm}
\begin{circuitikz}[scale=0.7, transform shape]
  \draw (0,0) -- (0,3) to[R, l=$R_1$] (2.5,3) to[R, l=$R_3$, i=$I_3''$] (5,3);
  \draw (2.5,3) to[R, l_=$R_2$, v^=$U_A''$, *-*] (2.5,0);
  \draw (5,3) to[R, l_=$R_4$, v^=$U_B''$, *-*] (5,0);
  \draw (7,0) to[I, i=$I_q$] (7,3) -- (5,3);
  \draw (0,0) -- (7,0);
  \node at (3.5,-0.7) {nur $I_q$ aktiv ($U$ = Kurzschluss)};
\end{circuitikz}
\end{center}
\schritt{b) Nur $U$ aktiv}
\leitfrage{Welchen Widerstand „sieht“ der Knoten A nach rechts unten?}
$R_3$ und $R_4$ liegen in Reihe ($\SI{30}{\ohm}$), parallel dazu $R_2$: $\SI{20}{\ohm}\parallel\SI{30}{\ohm} = \SI{12}{\ohm}$. Spannungsteiler mit $R_1$:
\[
  U_A' = \SI{16}{V}\cdot\frac{\SI{12}{\ohm}}{\SI{20}{\ohm}+\SI{12}{\ohm}} = \SI{6}{V}, \qquad
  I_3' = \frac{U_A'}{R_3+R_4} = \frac{\SI{6}{V}}{\SI{30}{\ohm}} = \SI{0.2}{A}, \qquad
  U_B' = R_4\cdot I_3' = \SI{4}{V}
\]
\schritt{c) Nur $I_q$ aktiv}
Jetzt liegen $R_1$ und $R_2$ parallel ($\SI{10}{\ohm}$), in Reihe dazu $R_3$ ($\SI{20}{\ohm}$), und das parallel zu $R_4$:
\[
  R_B = \SI{20}{\ohm}\parallel\SI{20}{\ohm} = \SI{10}{\ohm} \quad\Rightarrow\quad U_B'' = I_q\cdot R_B = \SI{10}{V}
\]
Strom durch den linken Zweig (von B nach A, also \emph{gegen} den Pfeil $I_3$):
\[
  I_3'' = -\frac{U_B''}{R_3 + R_1\parallel R_2} = -\frac{\SI{10}{V}}{\SI{20}{\ohm}} = -\SI{0.5}{A}, \qquad
  U_A'' = \SI{0.5}{A}\cdot\SI{10}{\ohm} = \SI{5}{V}
\]
\schritt{d) Überlagerung}
\[
  U_A = \SI{6}{V} + \SI{5}{V} = \ergebnis{\SI{11}{V}}, \qquad
  U_B = \SI{4}{V} + \SI{10}{V} = \ergebnis{\SI{14}{V}}, \qquad
  I_3 = \SI{0.2}{A} - \SI{0.5}{A} = \ergebnis{-\SI{0.3}{A}}
\]
Das negative Vorzeichen bedeutet: Der Strom fließt tatsächlich von B nach A. Das passt zu $U_B > U_A$.
Kontrolle mit dem Ohmschen Gesetz: $I_3 = (U_A - U_B)/R_3 = \SI{-3}{V}/\SI{10}{\ohm} = \SI{-0.3}{A}$ \checkmark
\schritt{e) Leistung der Stromquelle}
An der Stromquelle liegt $U_B$ an; Strom- und Spannungspfeil (gedacht von oben nach unten) bilden ein Erzeugerzählpfeilsystem:
\[
  P_{I_q} = U_B\cdot I_q = \SI{14}{V}\cdot\SI{1}{A} = \ergebnis{\SI{14}{W}} \ \text{(abgegeben)}
\]
Die Spannung an der Stromquelle stellt sich erst durch die \emph{gesamte} Schaltung ein -- auch hier darf man nicht mit Teilergebnissen rechnen.
\typfehler{Die Stromquelle durch einen Kurzschluss ersetzen -- dann wäre $R_4$ kurzgeschlossen und $U_B' = 0$. Merke: Spannungsquelle $\to$ Kurzschluss, Stromquelle $\to$ Leerlauf.}
\end{loesung}

% -------------------------------------------------------------------
\begin{aufgabe}{Drei Quellen an einem Knoten}{17}
\begin{center}
\begin{circuitikz}
  \draw (0,0) to[V, v<=$U_1$] (0,2) to[R, l=$R_1$] (0,4) -- (7.5,4);
  \draw (2.5,0) to[V, v=$U_2$] (2.5,2) to[R, l=$R_2$] (2.5,4);
  \draw (5,4) to[R, l_=$R_3$, v^=$U_A$, i>_=$I_3$, *-*] (5,0);
  \draw (7.5,0) to[I, i=$I_q$] (7.5,4);
  \draw (0,0) -- (7.5,0);
  \fill (2.5,4) circle (1.5pt) (2.5,0) circle (1.5pt);
\end{circuitikz}
\end{center}
\gegeben{$U_1 = \SI{12}{V}$, $U_2 = \SI{6}{V}$, $I_q = \SI{2}{A}$, $R_1 = \SI{6}{\ohm}$, $R_2 = \SI{3}{\ohm}$, $R_3 = \SI{2}{\ohm}$}
\begin{teilaufgaben}
  \item Achten Sie auf die Polung von $U_2$! Welches Vorzeichen erwarten Sie für den Beitrag von $U_2$ zu $U_A$?
  \item Berechnen Sie die drei Teilbeiträge $U_A'$ (nur $U_1$), $U_A''$ (nur $U_2$) und $U_A'''$ (nur $I_q$).
  \item Geben Sie $U_A$ und $I_3$ an.
  \item Auf welchen Wert muss $U_1$ erhöht werden, damit $I_3 = \SI{2}{A}$ wird? Nutzen Sie die Linearität aus, ohne neu zu rechnen.
\end{teilaufgaben}
\end{aufgabe}

\begin{loesung}
\schritt{a) Polung von $U_2$}
Der Spannungspfeil von $U_2$ zeigt nach oben, der Pluspol liegt also \emph{unten}. $U_2$ zieht den Knoten A unter das Bezugspotential -- der Beitrag $U_A''$ ist negativ.
\schritt{b) Drei Teilschaltungen}
\begin{center}
\begin{circuitikz}[scale=0.7, transform shape]
  \draw (0,0) to[V, v<=$U_1$] (0,2) to[R, l=$R_1$] (0,4) -- (4,4);
  \draw (2,0) node[circ]{} -- (2,2) to[R, l=$R_2$, -*] (2,4);
  \draw (4,4) to[R, l_=$R_3$, v^=$U_A'$] (4,0) -- (0,0);
  \node at (2,-0.8) {\Large nur $U_1$};
\end{circuitikz}
\hfill
\begin{circuitikz}[scale=0.7, transform shape]
  \draw (0,0) -- (0,2) to[R, l=$R_1$] (0,4) -- (4,4);
  \draw (2,0) to[V, v=$U_2$] (2,2) to[R, l=$R_2$, -*] (2,4);
  \draw (4,4) to[R, l_=$R_3$, v^=$U_A''$] (4,0) -- (0,0);
  \fill (2,0) circle (2pt);
  \node at (2,-0.8) {\Large nur $U_2$};
\end{circuitikz}
\hfill
\begin{circuitikz}[scale=0.7, transform shape]
  \draw (0,0) -- (0,2) to[R, l=$R_1$] (0,4) -- (6,4);
  \draw (2,0) node[circ]{} -- (2,2) to[R, l=$R_2$, -*] (2,4);
  \draw (4,4) to[R, l_=$R_3$, v^=$U_A'''$, *-*] (4,0);
  \draw (6,0) to[I, i=$I_q$] (6,4);
  \draw (0,0) -- (6,0);
  \node at (3,-0.8) {\Large nur $I_q$};
\end{circuitikz}
\end{center}
Nur $U_1$: $R_2\parallel R_3 = \SI{3}{\ohm}\cdot\SI{2}{\ohm}/\SI{5}{\ohm} = \SI{1.2}{\ohm}$
\[
  U_A' = U_1\cdot\frac{\SI{1.2}{\ohm}}{\SI{6}{\ohm}+\SI{1.2}{\ohm}} = \SI{12}{V}\cdot\frac{1}{6} = \SI{2}{V}
\]
Nur $U_2$: $R_1\parallel R_3 = \SI{6}{\ohm}\cdot\SI{2}{\ohm}/\SI{8}{\ohm} = \SI{1.5}{\ohm}$, Vorzeichen negativ (siehe a):
\[
  U_A'' = -U_2\cdot\frac{\SI{1.5}{\ohm}}{\SI{3}{\ohm}+\SI{1.5}{\ohm}} = -\SI{6}{V}\cdot\frac{1}{3} = -\SI{2}{V}
\]
Nur $I_q$: Alle drei Widerstände liegen parallel, $R_1\parallel R_2\parallel R_3$: $\frac{1}{6}+\frac{1}{3}+\frac{1}{2} = \SI{1}{S}$, also $\SI{1}{\ohm}$:
\[
  U_A''' = I_q\cdot\SI{1}{\ohm} = \SI{2}{V}
\]
\schritt{c) Überlagerung}
\[
  U_A = \SI{2}{V} - \SI{2}{V} + \SI{2}{V} = \ergebnis{\SI{2}{V}}, \qquad I_3 = \frac{U_A}{R_3} = \ergebnis{\SI{1}{A}}
\]
\schritt{d) Linearität ausnutzen}
\leitfrage{Wie hängt der Beitrag $U_A'$ von $U_1$ ab?}
$I_3 = \SI{2}{A}$ verlangt $U_A = \SI{4}{V}$. Die Beiträge von $U_2$ und $I_q$ ändern sich nicht (zusammen $\SI{0}{V}$), also muss $U_A' = \SI{4}{V}$ werden. Da $U_A'$ proportional zu $U_1$ ist ($U_A' = U_1/6$), muss $U_1$ verdoppelt werden:
\[
  U_1 = \ergebnis{\SI{24}{V}}
\]
\typfehler{Bei „nur $I_q$“ den Zweig mit $R_1$ oder $R_2$ vergessen: Die kurzgeschlossenen Spannungsquellen verbinden diese Widerstände mit dem unteren Knoten, sie bleiben also Teil der Schaltung. Bei d) die Linearität auf das Gesamtergebnis anwenden ($U_1$ verdoppeln $\Rightarrow$ $U_A$ verdoppelt) -- das gilt nur, wenn \emph{alle} Quellen gemeinsam skaliert werden.}
\end{loesung}

% -------------------------------------------------------------------
\begin{aufgabe}{Verständnisfragen}{8}
Sind die folgenden Aussagen richtig oder falsch? Begründen Sie kurz.
\begin{teilaufgaben}
  \item Beim Überlagerungssatz werden nicht betrachtete Spannungsquellen durch einen Leerlauf ersetzt.
  \item Der Innenwiderstand einer \emph{realen} Quelle bleibt beim Deaktivieren in der Schaltung.
  \item Der Überlagerungssatz gilt auch für Schaltungen mit Glühlampen, deren Widerstand von der Temperatur abhängt.
  \item Hat eine Schaltung drei Quellen, muss man drei Teilschaltungen berechnen.
\end{teilaufgaben}
\end{aufgabe}

\begin{loesung}
\schritt{a) Falsch}
Eine deaktivierte Spannungsquelle hat $U = 0$ -- das ist ein \emph{Kurzschluss}. Durch einen Leerlauf ersetzt man Stromquellen ($I = 0$).
\schritt{b) Richtig}
Deaktiviert wird nur die \emph{ideale} Quelle. Der Innenwiderstand ist ein gewöhnlicher Widerstand und bleibt in der Schaltung.
\schritt{c) Falsch}
Der Überlagerungssatz setzt eine \emph{lineare} Schaltung voraus. Ein temperaturabhängiger Widerstand ändert seinen Wert mit dem Strom, die Teilergebnisse dürfen dann nicht addiert werden.
\schritt{d) Richtig (in der Regel)}
Jede Quelle ist einmal allein aktiv. Man darf aber auch Gruppen bilden, z.\,B.\ zwei Quellen gemeinsam aktiv lassen und die dritte deaktivieren -- dann genügen zwei Teilschaltungen.
\typfehler{Kurzschluss und Leerlauf vertauschen ist der häufigste Fehler beim Überlagerungssatz. Merkhilfe: Die deaktivierte Quelle liefert „nichts“ -- eine Spannungsquelle also $\SI{0}{V}$ (Draht), eine Stromquelle $\SI{0}{A}$ (Unterbrechung).}
\end{loesung}

\end{document}
